\documentclass[man,11pt,a4paper]{apa7}
% APA 7 professional-manuscript layout, set in a sans-serif (Helvetica, metrically
% identical to Arial). Converted from the Word original; wording unchanged.
\usepackage[T1]{fontenc}
\usepackage[utf8]{inputenc}
\usepackage[scaled]{helvet}
\renewcommand{\familydefault}{\sfdefault}
\usepackage{url}
\urlstyle{same}

\title{Evaluating Extreme Technological Risk: A Social Contract Based Approach}
\shorttitle{EVALUATING EXTREME TECHNOLOGICAL RISK}
\authorsnames{SJ Beard, Patrick Kaczmarek}
\authorsaffiliations{{Centre for the Study of Existential Risk, University of Cambridge}}
\abstract{Philosophers engaging with the most extreme forms of technological risk, and in particular with existential risk to humanity as a whole, tend to assume a consequentialist perspective on the grounds that it alone can take account of the harm such technology might cause to possible future people. We believe that this is a mistake, and that other ethical theories are no less able to account for the full badness of these risks. In this paper we consider extreme technological risk from a social contract perspective, and argue that we may have obligations to past, present, and future people, as well as to humanity as a whole, to prioritize the mitigation of extreme technological risk and advance the interests of future people.}
\authornote{A version of this paper was published as Chapter 4, Evaluating Extreme Technological Risk: A Social Contract Based Approach, in Catherine Rhodes ed (2024) \textit{Managing Extreme Technological Risk}. World Scientific (pp. 109--126).}

\begin{document}
\maketitle
% SJ's choice (not APA): single-spaced footnotes
\makeatletter
\renewcommand\@makefntext[1]{\linespread{1}\selectfont\raggedright\textsuperscript{\@thefnmark}~#1}
\makeatother

\section{Introduction}

Extreme technological risks necessarily involve significant intergenerational considerations. On the one hand, the development of risky technologies, like biotechnology, nanotechnology and artificial intelligence (AI), is often justified by their ability to alter the trajectory of humanity's future towards what is maximally beneficial and safe for us. On the other hand, these technologies carry the potential to cause global catastrophes, such as human extinction, the collapse of human civilization or a permanent shift in humanity's future trajectory that would be harmful and/or dangerous. While these changes would most likely carry significant impact on the present generation, that is only a small part of their effects and we cannot truly evaluate them without taking the intergenerational perspective (Parfit, 1984).

The standard means for undertaking such an evaluation is by impartially weighing up costs and benefits from a consequentialist moral standpoint, such as utilitarianism. This has been, implicitly or explicitly, the position taken by the majority of researchers in the field of Existential Risk Studies, and while there are undoubtedly some important philosophical considerations that need to be addressed when carrying out such evaluations, the results have been generally serviceable and well founded (Millett \& Snyder-Beattie, 2017; Wolff, 2021). Nevertheless, it would be bad if this was the only way in which we could capture the moral importance of these aspects of extreme technological risk. For one thing, consequentialist approaches to moral theory in general, and cost-benefit analysis in particular, are controversial and many people do not accept them. While there is general academic acceptance of these views across, say, most leading economics departments, many anglophone philosophy departments, much of contemporary `start-up culture' and regulatory institutions in the UK and USA, there is profound scepticism and even hostility towards them among most leading public health departments, many non-anglophone philosophy departments, most of contemporary `activist culture' and regulatory institutions in continental Europe. Another problem is that each of us should be humble enough to recognize that an attitude of moral uncertainty is epistemically appropriate, and that even if we were to broadly agree with consequentialism (which the authors do), it may still turn out to be incorrect and superseded by other moral theories that could be very different in both substance and spirit.

For these reasons, we consider it prudent for philosophers with an interest in extreme technological risk to constructively engage with other moral perspectives and develop evaluative approaches that are compatible with these. In service of that goal, we here summarize one such approach based upon the social contract tradition.

\section{The Intergenerational Nature of the Social Contract}

The social contract tradition holds, roughly, that the content of morality derives from agreement between all those in the moral domain (Ashford \& Mulgan, 2018). If the `social contract' were understood as an \textit{actual} agreement made between various parties, and as such required that only those parties who could presently submit to, or negotiate, an agreement can be understood as parties to it, then this would seem to rule out any effects of our actions on merely possible future people as an irrelevant moral consideration.

It is true that the pioneers of this tradition, such as Thomas Hobbes (1651) and Jean-Jacques Rousseau (1762), did appear to conceive of the social contract as some form of actual agreement that people make, so as to either save themselves from anarchy (according to Hobbes) or trade their natural freedoms for a better quality of life (according to Rousseau). However, from early on, people saw that even \textit{if} this made sense at a moral level it was clearly a misrepresentation of how we actually live our lives. As David Hume, writing about the social contract in 1748, argued:

\begin{quote}
Were you to preach ... that political connections are founded altogether on voluntary consent or a mutual promise, the magistrate would soon imprison you ... if your friends did not before shut you up as delirious for advancing such absurdities. (Hume, 1987, p. 470)
\end{quote}

Since then, nearly everyone who has engaged with social contract theory has claimed, not that this is an actual contract to which everyone consents, but rather a hypothetical contract that people either \textit{would} agree to under certain conditions, \textit{could} agree to without violating certain stipulations or \textit{ought} to agree to if they thought about it in a certain way. Understood in this sense, the possibility of a social contract that extends beyond those presently alive becomes much less problematic, and indeed such a contract was proposed by Edmund Burke as early as 1790.

``Society is indeed a contract,'' Burke maintains; however, it is different in both form and scope to the kinds of contract we are used to dealing with. While such ``[s]ubordinate contracts for objects of mere occasional interest may be dissolved at pleasure'' society ``is to be looked on with other reverence, because it is not a partnership in things subservient only to the gross animal existence of a temporary and perishable nature'' (Burke, 1999, pp. 192--193). Rather:

\begin{quote}
It is a partnership in all science; a partnership in all art; a partnership in every virtue and in all perfection. As the ends of such a partnership cannot be obtained in many generations, it becomes a partnership not only between those who are living, but between those who are living, those who are dead, and those who are to be born. (Burke, 1999, pp. 192--193)
\end{quote}

In other words, to be alive is not to be engaged in some solitary pursuit to meet your own present needs, but rather to engage in cooperative and perpetual projects for the betterment of all. We ourselves, the present generation, have benefited enormously from previous generations who have engaged in such projects, while future generations are relying upon us to do the same. What is more, this process of working for the good of posterity is not simply some happy accident of human nature (akin to Adam Smith's `invisible hand'), rather we choose to do this because it matters to us. As Samuel Scheffler (2010, p. 269) puts it, ``what would it mean to value things but, in general, to see no reason of any kind to sustain them or retain them or preserve them or extend them into the future?''

This notion of humanity as involved in some perpetual collective partnership to realize `Burkean projects', such as science, art, virtue and perfection, for the benefit of future generations is quite common amongst futurists, as is the notion that either causing or allowing a catastrophe to threaten these projects would betray the trust that previous generations have placed in us to complete what they began. For instance, Wendell Bell (1993) and Richard Slaughter (1994) both support the notion that present generations have an obligation to see that humanity's `important business' is not left unfinished. While Bruce Tonn, combining the work of these two with his own Rawlsian analysis and some basic assumptions about the future trajectory of humanity, argued that our obligations under such a partnership imply that ``the acceptable probability of human extinction for the imaginable future is $2 \times 10^{-20}$'' (2009, p. 432). Nick Bostrom has similarly claimed that ``we might also have custodial duties to preserve the inheritance of humanity passed on to us by our ancestors and convey it safely to our descendants. We do not want to be the failing link in the chain of generations'' (2013, p. 23). While, most recently, Leopold Aschenbrenner has developed a version of these ideas he calls Burkean Longtermism (2021).

Yet, such lines of reasoning have received little serious attention within analytic moral philosophy; although there are clear parallels in other, more communitarian, traditions to intergenerational ethics including, but not limited to, those explored in contemporary African philosophy (Behrens, 2012). In what follows, we offer a preliminary analysis of how such a theory might be developed by summarizing three key lines of argument concerning the nature of our duties to the past, the future and humanity as a whole.\footnote{These rehearse arguments that we have made elsewhere, and we are grateful to the editors of those journals, \textit{Utilitas} and \textit{Argumenta}, for permission to reprint them.} These arguments are developed to be compatible with a particular version of the social contract tradition associated with the work of Tim Scanlon (1998). While not necessarily reflecting all of social contract theory as a whole we chose this form of contractualism for three reasons. Firstly, it is the most commonly cited in contemporary analytical moral philosophy. Secondly, it is the form of contractualism which what little literature has been produced on these issues tends to engage with. Thirdly, it allows for a greater degree of formalization than most versions of contractualism, and while we offer only the bare minimum of formality, we hope eventually to develop a theory that can be put to work on technical issues requiring such formality, such as risk assessment and population ethics.

\section{Our Obligations to the Past}

Consider:

\begin{quote}
Through no fault or choice of his own, Jeff is very sick. He desperately needs a liver transplant. Although he is not obliged to do so, a stranger called Michael gives Jeff part of his liver at the cost of reducing his own lifespan by ten years. After the procedure, Jeff chooses to drink heavily, and he dies from cirrhosis shortly afterwards.
\end{quote}

In considering a case like this, it is our view that Michael does the right thing by Jeff in donating part of his liver but that, by being reckless, Jeff does the wrong thing, not only by himself but also by Michael, and this makes the state of affairs worse than it would otherwise have been.

Let's start by distinguishing two features of this case that contribute to the moral status of Michael's sacrifice in donating part of his liver to Jeff. Firstly, there is the question of whether Michael acted appropriately as a moral agent. The answer to this question doesn't depend on what Jeff does next. All that matters is whether Michael himself is sensitive to the moral reasons involved in this decision (including the expected benefit to Jeff at the time of transplant, the expected cost to Michael and any further considerations, such as special duties Michael may have to Jeff). Secondly, we can ask whether other people in this case are sensitive to the moral reasons that flow from Michael's decision, and in particular the fact that, by deciding as he has, Michael has been harmed.

Despite the fact that it was morally justified, shortening Michael's life is a bad thing that, all else being equal, we would remove from the world if we could. Jeff's actions post-surgery appear to make this bad thing even worse, not because they make Michael's life any worse, but because they fail to realize the full extent of the benefits made possible by Michael's sacrifice. When choosing to donate part of his liver, Michael did the right thing because he judged that Jeff could expect to get much more than the ten years of additional life that Michael himself had sacrificed. He trusted Jeff to realize this full benefit from his sacrifice, making it, in other words, worthwhile. However, Jeff's actions post-surgery meant that he actually benefited by a lot less than Michael had lost. Jeff thus betrayed Michael's trust, and thereby wronged him.

What is true in this rather simple case also applies, we would argue, in cases at a much grander scale. Although our actions today can neither prevent nor mitigate the harms that were involved in past sacrifices, they can still do right or wrong by those who made those sacrifices by supporting or frustrating the Burkean projects they were engaged in, and thus making their losses more or less worthwhile. Our ancestors laboured hard under conditions we would find intolerable to secure benefits for their descendants, fought to end slavery and promote human rights, participated in the creation of feats of engineering, artistry, science and technological development whose benefits they would never see, and otherwise made numerous sacrifices for our sake. The worthwhileness of such sacrifices could well alter the moral reasons we have for acting to varying degrees and, when contravened in the absence of adequate justification, can indicate a lack of concern for the interests of others. Let us call this the Moral Worthwhileness Intuition (hereafter abbreviated `MWI').

To accept MWI is to say that we have a \textit{pro tanto} reason to promote the interest of future generations as Burke suggested, given the many sacrifices that our forebears have made for their sake. Even if it would not wrong any future people if ours was to be the last generation of humans, it still seems to wrong those who had anticipated and made grave sacrifices for the benefit of future generations and the long-term future of our species. Only if future generations actually exist, and enjoy the fruits of these Burkean projects, can we make those sacrifices as worthwhile as possible.\footnote{In a reply to our arguments Finneron-Burns (2022) has suggested that we cannot make sense of this claim unless we advocate for the existence of `existential benefits', i.e. unless we believe that these projects will actually benefit future people who may not even exist without them. While we do believe in existential benefits, we do not think this is necessary for our claims. One can be inspired to make sacrifices towards achieving Burkean projects merely because one recognizes that these would be enjoyed by future people were they to exist, without saying that their enjoyment benefits such people. Others might then be obliged to make these sacrifices worthwhile out of a duty to past sacrificers, even if they do not believe in existential benefits.}

\subsection{The Value of Worthwhileness}

Yet, merely accepting MWI still leaves important questions to be answered about the nature of our obligations to the past, and whether these give us sufficiently strong reason to both work for the benefit of posterity and to prevent catastrophes that would destroy humanity's potential.

One can appreciate the general form that such a theory might take by considering the fact that one who has made a sacrifice would reasonably object to another person acting in such a way as to make their sacrifice worthless. This is because making this sacrifice was morally supererogatory and done purely because of the benefits it would bring to others. Actually realizing these benefits, however, was not fully in the gift of the sacrificer and is thus left as \textit{unfinished business} for others to complete: including both the person benefited and potentially other intermediaries between the sacrificer and the beneficiary. This unfinished business takes the form of additional obligations to act in ways that might otherwise have been beyond the call of duty.

To put this more formally, on the \textit{Unfinished Business Account}:

\begin{quote}
All else being equal, if actor \textit{p} reasonably judges that performing a supererogatory act \textit{$\phi$} at great sacrifice to herself will enable beneficiary \textit{q} to achieve some good, then failure to promote the good made possible by \textit{$\phi$} wrongs \textit{p, even if p is no longer alive}.\footnote{In an earlier version of this argument we instead said ``will enable beneficiary \textit{q} to achieve a greater good''. We have since reflected however that it does not seem right to invalidate the Unfinished Business Account merely because the original sacrificer chose, imprudently, to sacrifice more than could be gained from their sacrifice. Such acts are imprudent to be sure, and we may have reasons to discourage them, but we may still have duties to render them as worthwhile as we can anyway.}
\end{quote}

Of course, this does not imply that it is always impermissible to wrong someone in this way: I cannot oblige another person to do my bidding simply by making some mild sacrifice whose benefits would depend greatly on what she then chose to do. We only have obligations to the past when the reasons that flow from their sacrifices are sufficiently strong and decisive, and this will depend upon at least three things: (1) how worthwhile the sacrifice would be with or without our intervention, (2) the cost to us of intervening and (3) the degree of fit between such an obligation and our long-term interests and conception of the good life (that is, our disposition to bear such burdens even if they were not obligatory). At the limit it is clearly not possible to oblige another person to accept a greater sacrifice than that which one originally accepted or to accept a sacrifice that, when combined with our own, would yield a total cost that was greater than the total benefit we could together have achieved. While Michael could reasonably expect Jeff to refrain from drinking himself to death as a result of sacrificing part of his liver, he could not expect Jeff to undergo torture, even if this would ultimately produce some even greater good.

More substantially, it seems only reasonable that if our obligations to the past are to play a decisive role in our moral choices, then they should be at least broadly in line with our long-term interests and consistent with our conception of the good life, even if the required action is not what we ourselves would have chosen. For instance, we can easily imagine that Jeff looking after his liver is consistent with his own conception of the good life, even if he would actually much rather stare down the bottom of a bottle, but were we to require him to convert to Catholicism and eat nothing but ramen noodles this would not be. At the grander scale this implies that, while we may be bound to complete projects of genuinely shared human value and to promote the general interests of future people, we have no obligation to realize the mere whims or peculiar interests of our forebears concerning the future.

One may wonder whether we really can have obligations to those who have made sacrifices even once they are dead. There is nothing either new or radical in this view (see, e.g., Marquis, 1985); however, there are some philosophers who claim that this is not possible because, for instance, ``an act cannot be wrong if it would be worse for no one'' (Parfit, 2017, p. 118). We can never wrong our forebears, according to such views, because nothing we now do can make them better or worse off; after all, they are dead. However, it is unclear just how substantial this objection proves, given that, even if we don't think they are directly harmed, most of us do set great store by respecting the past and those who lived in it.

Nevertheless, this deserves some exploration. Here, we assume the truth of Carl Wellman's notion of \textit{surviving duties}; as he writes, ``[d]eath destroys a person's capacity for agency and thereby destroys an essential qualification for being a moral right-holder. At the same time, death sometimes leaves the moral reasons intact so that, together with the qualifications of the surviving duty-bearer, these can ground a surviving duty'' (1995, pp. 156--157). Our claim is that, so far as we have a duty to realize the value of others' sacrifices to realize the grand projects of humanity, these duties do not perish with the death of the person who performed this sacrifice; nor does their strength depreciate.\footnote{For a selection of objections to this argument and our responses to them see Kaczmarek and Beard (2020).}

\section{Our Obligations to the Future}

Still, as Burke says, our obligations under the social contract do not only extend to those who are living and those who are dead but also `those who are to be born'. Here, the potential for our choices to make things worse is not in doubt. However, there are still those who object that, from the point of view of the social contract, these negative impacts are not morally salient.

The most common defence of this claim stems from the Scanlonian version of contractualism according to which, ``an act is wrong if its performance under the circumstances would be disallowed by any set of principles for the general regulation of behaviour that no one could reasonably reject as a basis for informed, unforced general agreement'' (Scanlon, 1998, p. 153). This is to say, a principle is morally sound just in case it is justifiable to each person, where such justification is attentive to the interests, complaints and other claims of the individuals in question.

To define this view more precisely, Scanlon places two restrictions on the reasons that can be offered for rejecting a moral principle. Firstly, there is what Derek Parfit has called \textit{The Impersonalist Restriction:}

\begin{quote}
Impersonal values are not themselves grounds for reasonable rejection. (Scanlon, 1998, p. 222; cf. Parfit, 2011, p. 214)
\end{quote}

This means that if it's wrong to fail to bring future people into existence, then this must be because of the value of this existence to these people. It cannot simply be a recognition that these people's lives would be good from some third-party point of view: as Henry Sidgwick memorably put it, `the point of view of the universe' (Sidgwick, 1907, p. 382).

Moreover, there is what Parfit dubs \textit{The Individualist Restriction:}

\begin{quote}
The justifiability of a moral principle depends only on individuals' reasons for objecting to that principle and alternatives to it. (Scanlon, 1998, p. 229; cf. Parfit, 2011, p. 193)
\end{quote}

This means that the value of future people's lives to them must rest in the fact that this promotes some interest that these people actually have, and that we are failing to take into account in the outcome in which these individuals never exist.

\subsection{Finneron-Burns' Argument}

In a recent paper (2017), Elizabeth Finneron-Burns has developed a series of arguments based on these restrictions, which aim to show that the potential loss of future generations, including from possible technological catastrophes, should not (that is, non-derivatively) factor into our choices. In particular, she argues that it would be morally irrelevant that our actions might lead to many people never coming into existence because there would be no one \textit{there} who has an interest in being brought into existence, and whom we failed to attend to. To be sure, she then goes on to argue that causing such a disaster would still be wrong, but only due to its impact on those already living.

Finneron-Burns states her principal argument against the need to take account of the impact of catastrophes on merely possible people as follows:

\begin{quote}
[W]e can only wrong someone who did, does or will actually exist because wronging involves failing to take a person's interests into account. When considering the permissibility of a principle allowing us not to create Person X, we cannot take X's interest in being created into account because X will not exist if we follow the principle. By considering the standpoint of a person in our deliberations, we consider the burdens they will have to bear as a result of the principle. In this case, there is no one who will bear any burdens since if the principle is followed (that is if we do not create X), X will not exist to bear any burdens. So, only people who do/will actually exist can bear the brunt of a principle, and therefore occupy a standpoint that is owed justification. (Finneron-Burns, 2017, p. 331)
\end{quote}

According to this argument, we can only wrong people by failing to take their interests into account (i.e., the individualist restriction), and people's interests only extend to not being burdened. Since merely possible people do not have interests, as people who do not exist do not bear any burden by virtue of their non-existence, we, therefore, cannot wrong them.

\subsubsection{How We Might Still Have Obligations to Future People}

One reply to this argument comes from Rahul Kumar (2018). Kumar agrees with Finneron-Burns that we do not \textit{harm} future people by causing them not to exist. However, he argues that a contractualist can still object to causing or allowing catastrophes that would cause future people not to exist, even though this does not constitute a harm to them. There are cases, he argues, in which we act wrongly by failing to take a person's interests into account even though we do not, in fact, harm or burden them in any way. One example in which we can be said to do this is when we act negligently, or even viciously, towards a person but fail to do any harm through sheer dumb luck. If this is possible, the argument goes, then why is it not equally possible for us to fail to take somebody's interests into account in cases in which that person is not burdened because we cause them not to exist?

In an earlier paper (2016), Finneron-Burns objects to this line of argument on the grounds that, while it is possible to wrong a person without harming them, this can only be achieved if there is a relevant `standpoint', or cluster of personal characteristics and interests that an individual can occupy, whose interests are not being taken into account. There are, she admits, a great variety of such standpoints that may be morally salient. Any of these could give us reasons for rejecting a moral principle on the grounds that it does wrong by a person who might occupy this standpoint, even though it does not actually harm anyone because nobody, in the end, actually occupies it. However, she goes on to claim that this approach cannot be applied to the interests of merely possible people because standpoints can only describe positions that could be occupied by actual people and existence is not an interest that can attach to a standpoint (Finneron-Burns, 2016, pp. 1156ff.).

According to another reply, it is claimed that we fail to take individuals' interest into account in cases when they do not exist because non-existence is equivalent in value to a person existing with a neutral life, that is a life that is neither good nor bad for them, so that they would be genuinely better off if they had received a good life instead. On this account, we can be said to burden merely possible future people because we have brought about a state, non-existence, that is worse for them than a possible happy life, which they might have had if we had acted otherwise (Arrhenius \& Rabinowicz, 2015; Fleurbaey \& Voorhoeve, 2015; cf. Pummer, 2019). This makes non-existence a real burden compared with the possibility of existence with a good life, and if we choose this option, then we can be said to harm in the same way as if we had simply made the life of an actual person much worse so that it were no longer good for them.

However, from a Scanlonian perspective this raises the troubling prospect of ethical paralysis. Since everybody who does not come into existence with a life worth living is, on this account, harmed in just the same way as a person forced to live a neutral life, then it seems that all non-existing people have strong reasons to reject any principle that does not bring them into existence (cf. Kavka, 1975). Since all, or virtually all, principles may have some effect upon the identity of those who exist in the future, this would mean that there is always someone whose interests we are failing to take account of when applying such principles, and hence no, or almost no, ethical principles will be left on the table. As Finneron-Burns puts it ``even basic, everyday activities would be morally wrong since they prevented certain people from coming to exist'' (Finneron-Burns, 2016, p. 1157).

We believe that there is another, better, reply to Finneron-Burns' argument that explains why we should care about the loss of future generations. According to this response, we do wrong by a merely possible person whom we denied existence, not because we bring about a state of affairs that has zero value for them, but because we do not bring about a state of affairs that would have had a positive value for them when we could have done so (without sacrificing anything of comparable value).

When we make choices that determine whether a possible future person will come into existence with a life worth living, we choose between an outcome that will be good for them and one that will not be good for them; even though, since they do not exist in this second outcome, it will not be worse for them either. To put it another way, on this proposed view, we can be said to fail to take a person's interests into account because we have failed to benefit them, even though we have not burdened them either. This kind of argument was first proposed by Krister Bykvist (2007). However, we focus here on the version offered in a posthumously published paper by Derek Parfit (2017).

Parfit labels the kind of position presented by Finneron-Burns the `No Complaints Claim', which he neatly summarized by the slogan ``[a]n act cannot be wrong unless there is or will be someone whom this act has wronged'' (Parfit, 2017, p. 136). He admits that this view has considerable psychological and moral force; however, Parfit is keen to point out that it only describes part of morality. In particular, it rests upon the idea that what we owe each other is exhausted by the reasons people give us not to act badly towards them, what Parfit refers to as `reasons of non-maleficence'. In Scanlon's terms, these are the reasons that a person might have for rejecting the principle on which we are acting.

In some of his earlier writings Parfit appeared to accept Scanlon's line of argument, going so far as to say that the wrongness of denying future people the benefit of existence cannot stem from the personal reasons that these people themselves would give us to bring them into existence (Parfit, 2011, p. 8). However, in what would be his final paper, Parfit moved beyond this position and argued that we could have personal reasons to bring people into existence. These would be reasons of beneficence to do what is good for a person and promote their `personal goodness'. For him, the difference between these reasons and reasons of `non-maleficence' seemed, at most, a matter of degree rather than of type, so that if a moral theory holds that we should respect reasons of non-maleficence it should not also claim that reasons of beneficence have no moral weight whatsoever (Parfit, 2017, pp. 136ff.).

While he discussed this possibility, Parfit did not offer any firm conclusions about it. However, his line of argument appears to us to suggest that Finneron-Burns' argument is far weaker than she claims. Specifically, things are hardly so cut and dried for the Scanlonian contractualist. They may (and if Parfit is right, should) accept a `Wide Person-Affecting Principle', according to which:

\begin{quote}
One of two outcomes would be in one way worse if this outcome would be less good for people, by benefiting people less than the other outcome would have benefited people. (Parfit, 2017, p. 129)
\end{quote}

Unlike the No Complaints Claim, Parfit's proposed principle does not restrict us to reasons of non-maleficence and instead considers the whole of morality, including how we might benefit people. Of course, there is no great difficulty in saying that a person whom we might bring into existence and who would have a good life could give us reasons to choose to bring them into existence, on the grounds of beneficence, because they would be glad to be alive in much the same way that a person whom we might bring into an awful, wretched existence provides us with a reason to choose not to bring them into existence, on the grounds of non-maleficence, because they would regret our actions. To put this more formally:

\begin{quote}
If person \textit{p} is caused to exist and to have a life that is worth living, then that is good for this person, giving him or her an `existential benefit'; all else being equal, denying this benefit is wrong, even though \textit{p} is not made worse off as a result.
\end{quote}

Finneron-Burns claims that ``wronging involves failing to take a person's interests into account'' (2017, p. 331); however, when arguing that we cannot wrong future generations she equates a person's interests with the avoidance of bearing burdens. Yet, many of us would say that avoiding burdens constitutes an incomplete (nay, impoverished) picture, and that our interests can equally be served by receiving benefits. When we fail to bring somebody into existence we can legitimately be said to have failed to take their interests into account, even if we didn't harm them, because we also didn't benefit them (by bringing them into existence) when we might easily have done so. It is in this respect that the interests of possible people count as reasons in favour of avoiding the destruction of the future, and not, as Finneron-Burns contends it would have to be, that we should militate against such threats only to the extent that doing so avoids harm. Since one of the outcomes that we might choose is good for these people, and hence serves their interests, failing to bring this outcome about would fail to take account of their interests and thus constitute a kind of mean-spiritedness, even though the actual outcome is not worse for them. Hence, we can be said to have done wrong by failing to act in the way that would have been good for these people, and thus denied them this existential benefit, whether or not the interests we fail to take into account can be said to actually exist in the outcome that we eventually bring about.

Of course, Finneron-Burns and Scanlon have their reasons for preferring to focus on moral burdens. One of these is surely that their brand of morality, in general, has a negative tone, since the aim of Scanlonian contractualism is to identify which principles are morally impermissible by working out which ones people could reasonably reject (a formulation that traces back to Immanuel Kant rather than the social contract tradition). Parfit himself has done much to criticize the motivation for this stern view, and to argue that it must be counterbalanced by a moral philosophy that considers which principles are required by virtue of being the ones that people could `rationally accept', although he also believed that these two approaches to morality ultimately arrive at the same result; or at least that, as he once put it, ``[they] are climbing the same mountain on different sides'' (Parfit, 2011, p. 373).

However, we do not think that the point we are concerned with here requires us to accept this further claim of Parfit's. Perhaps the theoretical equivalence between impermissibility, reasonable rejection and \textit{burdening} is of symbolic importance. But there is no reason why it would not be of equal moral importance to equate impermissibility, reasonable rejection and \textit{not benefiting}. Indeed, for the kinds of so-called `same population' cases with which Scanlonian contractualists are overwhelmingly concerned, the two would appear to amount to very much the same thing, and it is only when one considers how our actions might affect the size, or even existence, of future generations that we see any difference in their practical implications. To claim that in these latter cases one can only go along with Scanlon's original concept of moral burden, and that since nobody is burdened by their non-existence, causing them not to exist cannot be wrong, appears to us to be going too far. Why do we not owe some justification to these possible people for failing to do what would have been good for them, and why would that not provide an equal ground for wrongness via reasonable rejection and hence to including their interests in our conception of the social contract?\footnote{Finneron-Burns provides further supplementary arguments, which we respond to in Beard and Kaczmarek (2019).}

\section{Our Obligations to Humanity as a Whole}

Thus, it seems plausible to us that Burke had it right when he asked us to view the social contract as something that extends across generations, and that we can have obligations to both the past and the future that require us to seek to promote the interests of future people while preserving humanity from extreme technological risks. However, we may also have reasons that go beyond even this expansive contract.

There are many ways in which we can think about the value of the future. According to a distinction drawn by Christine Korsgaard it is possible to both value something for its own sake, or as an end in itself, which is its `final value,' and to value it for the sake of something else, or as a means to an end, which is its `instrumental value' (Korsgaard, 1983). Our arguments so far have taken the form of claiming that we should seek to protect the future of humanity because it is good for people, promoting the interests of future generations and making the sacrifices of past and present generations more worthwhile. However, some philosophers have also argued that humanity as a whole also possesses final value, and that this is something the social contract requires us to promote, if we can, as well.

Johann Frick in particular has done most to advance this idea. The argument, as Frick presents it (2017), begins with the observation that people commonly attribute final value to a range of phenomena, such as languages, species and cultures, that they see as worthy of protection for their own sake. This suggests ``that humanity too, with its unique capacities for complex language use and rational thought, its sensitivity to moral reasons, its ability to produce and appreciate art, music, and scientific knowledge, its sense of history, and so on, should be deemed to possess final value'' (Frick, 2017, p. 359). And, heck, why not?

As we noted earlier, our intergenerational social contract exists not only because we happen to rely upon previous generations for our lives and livelihoods while future generations are reliant upon us for the same. It also exists because working for posterity and the betterment of humankind simply is something that matters to people (including ourselves) and that gives meaning and purpose to our lives. Having this collective Burkean project cut short is something that we should all, plausibly, be deeply concerned about for its own sake, and not simply because we owe it to the interests of others, especially if we were to discover that humanity alone possesses these capacities, sensitivities and abilities. As Frick says, that ``would be very bad, indeed one of the worst things that could possibly happen'' (Frick, 2017, p. 344).

\section{Conclusion}

The idea of an intergenerational social contract encompassing not only presently existing people but also those who have come before and those who are yet to come has a long history and seems to be attractive to many who worry about extreme technological risk. It seems a shame therefore that it has received so little attention within analytical moral philosophy, and that what attention there has been is so often used to provide arguments against it. We have tried to address some of these objections and to offer positive arguments for why our obligations to the past, the future, and humanity as a whole, concerning the prevention of extreme technological risk, form important aspects of the social contract. In doing so we hope to have provided a foundation for new ways of evaluating extreme technological risk that would appeal to communities who are opposed to utilitarianism, whether out of instinct, ideology or conviction.

\section{References}
\begin{list}{}{\setlength{\leftmargin}{0.5in}\setlength{\itemindent}{-0.5in}\setlength{\itemsep}{0pt}\setlength{\parsep}{0pt}}\linespread{2}\selectfont
\item Arrhenius, G., \& Rabinowicz, W. (2015). The value of existence. In I. Hirose \& J. Olson (Eds.), \textit{The Oxford handbook of value theory} (pp. 424--443). Oxford University Press.
\item Aschenbrenner, L. (2021, July 26). Burkean longtermism. \textit{For Our Posterity}. \url{https://www.forourposterity.com/burkean-longtermism/}
\item Ashford, E., \& Mulgan, T. (2018). Contractualism. In E. N. Zalta (Ed.), \textit{The Stanford encyclopedia of philosophy} (Summer 2018 ed.). \url{https://plato.stanford.edu/archives/sum2018/entries/contractualism/}
\item Beard, S., \& Kaczmarek, P. (2019). On the wrongness of human extinction. \textit{Argumenta}, \textit{5}(1), 85--97.
\item Behrens, K. G. (2012). Moral obligations towards future generations in African thought. \textit{Journal of Global Ethics}, \textit{8}(2--3), 179--191.
\item Bell, W. (1993). Why should we care about future generations? In H. F. Didsbury, Jr. (Ed.), \textit{The years ahead} (pp. 25--41). World Future Society.
\item Bostrom, N. (2013). Existential risk prevention as global priority. \textit{Global Policy}, \textit{4}(1), 15--31.
\item Burke, E. (1999). \textit{Select works of Edmund Burke: Vol. 2. Reflections on the revolution in France}. Liberty Fund.
\item Bykvist, K. (2007). The benefits of coming into existence. \textit{Philosophical Studies}, \textit{135}(3), 335--362.
\item Finneron-Burns, E. (2016). Contractualism and the non-identity problem. \textit{Ethical Theory and Moral Practice}, \textit{19}, 1151--1163.
\item Finneron-Burns, E. (2017). What's wrong with human extinction? \textit{Canadian Journal of Philosophy}, \textit{47}(2--3), 327--343.
\item Finneron-Burns, E. (2022). Human extinction and moral worthwhileness. \textit{Utilitas}, \textit{34}(1), 105--112.
\item Fleurbaey, M., \& Voorhoeve, A. (2015). On the social and personal value of existence. In I. Hirose \& A. Reisner (Eds.), \textit{Weighing and reasoning: Themes from the philosophy of John Broome} (pp. 95--109). Oxford University Press.
\item Frick, J. (2017). On the survival of humanity. \textit{Canadian Journal of Philosophy}, \textit{47}(2--3), 344--367.
\item Hobbes, T. (1651). \textit{Leviathan or the matter, forme and power of a commonwealth ecclesiasticall and civil}.
\item Hume, D. (1987). \textit{Essays, moral, political, and literary}. Liberty Fund.
\item Kaczmarek, P., \& Beard, S. (2020). Human extinction and our obligations to the past. \textit{Utilitas}, \textit{32}(2), 199--208.
\item Kavka, G. (1975). Rawls on average and total utility. \textit{Philosophical Studies}, \textit{27}(4), 237--253.
\item Korsgaard, C. M. (1983). Two distinctions in goodness. \textit{The Philosophical Review}, \textit{92}(2), 169--195.
\item Kumar, R. (2018). Risking future generations. \textit{Ethical Theory and Moral Practice}, \textit{21}(2), 245--257.
\item Marquis, D. (1985). Harming the dead. \textit{Ethics}, \textit{96}(1), 159--161.
\item Millett, P., \& Snyder-Beattie, A. (2017). Existential risk and cost-effective biosecurity. \textit{Health Security}, \textit{15}(4), 373--383.
\item Parfit, D. (1984). \textit{Reasons and persons}. Oxford University Press.
\item Parfit, D. (2011). \textit{On what matters} (Vol. 2). Oxford University Press.
\item Parfit, D. (2017). Future people, the non-identity problem, and person-affecting principles. \textit{Philosophy \& Public Affairs}, \textit{45}(2), 118--157.
\item Pummer, T. (2019). The worseness of nonexistence. In E. Gamlund \& C. Solberg (Eds.), \textit{Saving people from the harm of death} (pp. 215--228). Oxford University Press.
\item Rousseau, J.-J. (1762). \textit{Du contrat social; ou Principes du droit politique}.
\item Scanlon, T. M. (1998). \textit{What we owe to each other}. Harvard University Press.
\item Scheffler, S. (2010). \textit{Equality and tradition}. Oxford University Press.
\item Sidgwick, H. (1907). \textit{The methods of ethics} (7th ed.). Macmillan.
\item Slaughter, R. A. (1994). Why we should care for future generations now. \textit{Futures}, \textit{26}(10), 1077--1085.
\item Tonn, B. E. (2009). Obligations to future generations and acceptable risks of human extinction. \textit{Futures}, \textit{41}(7), 427--435.
\item Wellman, C. (1995). \textit{Real rights}. Oxford University Press.
\item Wolff, J. (2021). Risk and the regulation of new technologies. In T. Matsuda, J. Wolff, \& T. Yanagawa (Eds.), \textit{Risks and regulation of new technologies} (pp. 3--18). Springer.
\end{list}

\end{document}
